\section*{Chapitre \ref{ch:b2:magmatic-differentiation} --- Différenciation magmatique}

\begin{solution}{exo:b2:magmatic-differentiation:1}
(a) Le liquidus de An$_{50}$ est \qty{1450}{\celsius}. Le premier cristal se place sur le
solidus à cette température, entre An$_{66.7}$ (1372) et An$_{83.3}$ (1465)~:
$66.7 + 16.6\times(1450 - 1372)/(1465 - 1372) = 66.7 + 13.9 = \qty{80.6}{\percent}$
d'anorthite. (b) À \qty{1394}{\celsius}, le liquide est sur le liquidus, en An$_{33.3}$~; les cristaux sur le solidus, $66.7 + 16.6\times22/93 =
\qty{70.6}{\percent}$. Par la règle des leviers, la fraction cristallisée vaut $(50 -
33.3)/(70.6 - 33.3) = 16.7/37.3 = 0.448$. (c) Le dernier liquide disparaît quand les
cristaux atteignent An$_{50}$, à la température du solidus \qty{1287}{\celsius}~;
le liquidus à cette température se situe entre An$_{11.1}$ (1265) et An$_{20}$ (1334)~:
$11.1 +
8.9\times22/69 = \qty{13.9}{\percent}$ d'anorthite.
\end{solution}

\begin{solution}{exo:b2:magmatic-differentiation:2}
$\bar D_{\ce{Sr}} = 0.15\times0.000063 + 0.35\times0.157 + 0.50\times2.7 = 0.00001 +
0.0550 + 1.35 = 1.40$~; $\bar D_{\ce{Zr}} = 0.0001 + 0.0683 + 0.0005 = 0.069$~;
$\bar D_{\ce{Rb}} = 0.00001 + 0.0012 + 0.0125 = 0.014$. Seul le strontium, retenu par le
plagioclase, est compatible ($\bar D > 1$)~; le zirconium et le rubidium sont incompatibles,
le rubidium davantage.
\end{solution}

\begin{solution}{exo:b2:magmatic-differentiation:3}
(a) $C_L/C_0 = 0.5^{D-1}$~: Sr $0.5^{0.40} = 0.758$ (il diminue)~; Zr $0.5^{-0.931} =
1.907$~; Rb $0.5^{-0.986} = 1.981$, proche de la limite 2. (b) Zr/Rb varie de
$0.5^{0.069 - 0.014} = 0.5^{0.055} = 0.963$, moins de \qty{4}{\percent} alors que chaque
élément a presque doublé~: comme l'indique le corollaire, le rapport de deux \omterm{def:b2:elements-in-the-earth:partition}{éléments incompatibles}
garde la valeur du \omterm{def:b2:magma-emplacement:magma}{magma} parent, et renseigne donc sur la source
plutôt que sur la cristallisation.
\end{solution}

\begin{solution}{exo:b2:magmatic-differentiation:4}
(a) \ce{SiO2} $+5.03$, \ce{MgO} $-1.95$, \ce{K2O} $+0.82$ (en pour cent), la potasse
multipliée par $1.67/0.85 = 1.96$. (b) Pour un élément parfaitement incompatible,
$C_L/C_0 = 1/F$, donc $F = 0.85/1.67 = 0.509$~: la moitié du \omterm{def:b2:magma-emplacement:magma}{magma} aurait dû cristalliser.
(c) Retirer \qty{12.5}{g} pour \qty{100}{g} laisse $F \approx 0.88$ et
n'élèverait la potasse que d'environ $1/0.88 = 1.14$~: la cristallisation seule ne peut pas
expliquer la hausse. La potasse est venue de l'extérieur, des xénolites granitiques, riches en
potasse (\qty{3.6}{\percent}), que le \omterm{def:b2:magma-emplacement:magma}{magma} a assimilés.
\end{solution}

\begin{solution}{exo:b2:magmatic-differentiation:5}
(a) $F = 0.8$, $D = 5$. Fractionnée~: $0.8^{4} = 0.410$~; à l'équilibre~:
$1/(5 + 0.8\times(1 - 5)) = 1/1.8 = 0.556$. (b) $F = 0.1$, $D = 0.01$. Fractionnée~:
$0.1^{-0.99} = 9.77$~; à l'équilibre~: $1/(0.01 + 0.1\times0.99) = 1/0.109 = 9.17$.
(c) Quand $F \to 0$, la loi de Rayleigh tend vers $0$ pour $D > 1$ et vers l'infini pour
$D < 1$, tandis que la loi d'équilibre tend vers $1/D$ (le dernier liquide en équilibre
avec des cristaux de composition globale). Seule la \omterm{def:b2:magmatic-differentiation:crystallisation}{cristallisation fractionnée} peut épuiser
un liquide en \omterm{def:b2:elements-in-the-earth:partition}{élément compatible}.
\end{solution}

\begin{solution}{exo:b2:magmatic-differentiation:6}
(a) Une formule unitaire contient 1.6 \ce{MgO}, 0.4 \ce{FeO} et un \ce{SiO2}~: $1.6\times
40.30 = 64.48$, $0.4\times71.84 = 28.74$, 60.08, soit au total \qty{153.30}{g}. Pourcentages
massiques \ce{MgO} \qty{42.06}{\percent}, \ce{FeO} \qty{18.75}{\percent},
\ce{SiO2} \qty{39.19}{\percent}. (b) Règle des leviers sur \ce{MgO}~: $x = (10.0 - 8.0)/(42.06
- 8.0) = 2.0/34.06 = 0.0587$, soit environ \qty{6}{\percent} d'olivine. (c) Silice du
liquide~: $(49.0 - 0.0587\times39.19)/(1 - 0.0587) = 46.70/0.9413 =
\qty{49.61}{\percent}$. Retirer de l'olivine, moins riche en silice que le basalte,
élève la silice du liquide, mais lentement.
\end{solution}

\begin{solution}{exo:b2:magmatic-differentiation:7}
(a) $F(0) = 0.768$~; $F(1000) = 0.1091 - 0.5906 + 0.7678 = 0.286$~; $F(2000) = 0.4364
- 1.1813 + 0.7678 = 0.023$. (b) Le discriminant vaut $(5.9064\times10^{-4})^2 -
4\times1.091\times10^{-7}\times0.7678 = 3.4885\times10^{-7} - 3.3507\times10^{-7} =
1.379\times10^{-8}$, de racine carrée $1.174\times10^{-4}$~; la plus petite racine est
$H = (5.9064 - 1.174)\times10^{-4}/(2.182\times10^{-7}) = \qty{2169}{m}$, l'horizon de
Sandwich. (c) $1 - 0.768 = 0.232$~: près d'un quart du \omterm{def:b2:magma-emplacement:magma}{magma} avait cristallisé
avant la formation de la base de la zone inférieure affleurante, dans la zone cachée
en dessous et dans la série de bordure le long des parois et du toit.
\end{solution}

\begin{solution}{exo:b2:magmatic-differentiation:8}
Le \omterm{def:b2:magmatic-differentiation:cumulate}{cumulat} contient $D\,C_L = D\,C_0F^{D-1}$, donc entre deux niveaux
$\ln(\mathrm{Ni}_2/\mathrm{Ni}_1) = (D - 1)\ln(F_2/F_1)$. (a) $D - 1 = \ln(54/177)/
\ln(0.29/0.70) = (-1.187)/(-0.881) = 1.35$, donc $D \approx 2.35$ (la régression sur les
66 \omterm{def:b2:magmatic-differentiation:cumulate}{cumulats} avec $F \ge 0.23$ donne 2.38). (b) $D - 1 = \ln(6.2/54)/\ln(0.17/0.29) = (-2.164)/(-0.534)
= 4.05$, donc $D \approx 5.05$. (c) Le $D$ global change quand l'assemblage cristallisant change
(un nouveau minéral qui fixe le nickel, ou un minéral qui disparaît), et les coefficients
minéral--liquide eux-mêmes augmentent quand le liquide se refroidit et change de composition,
comme cela a été mesuré pour l'olivine~; le liquide piégé dans les \omterm{def:b2:magmatic-differentiation:cumulate}{cumulats} brouille
l'estimation.
\end{solution}

\begin{solution}{exo:b2:magmatic-differentiation:9}
(a) $100 + 25.4 - 2.9 - 9.6 = \qty{112.9}{g}$. (b) En grammes de \ce{SiO2}~: $55.53 +
25.4\times0.7248 - 2.9\times0.3919 - 9.6\times0.5054 = 55.53 + 18.41 - 1.14 - 4.85 =
\qty{67.95}{g}$, donc $67.95/112.9 = \qty{60.19}{\percent}$ contre 60.16. \ce{MgO}~:
$4.66 + 0.20 - 1.22 = \qty{3.64}{g}$, soit \qty{3.22}{\percent} contre 3.19. \ce{K2O}~:
$0.88 + 0.92 = \qty{1.80}{g}$, soit \qty{1.60}{\percent} contre 1.75. (c) La silice et la magnésie, deux des trois
témoins, s'accordent à 0.03 près~; la potasse, non utilisée comme témoin, est trop basse de 0.15,
soit près d'un dixième de sa valeur~: le matériau assimilé avait plus de potasse que la xénolite moyenne,
ou un liquide riche en potasse issu de celle-ci a été prélevé de préférence.
\end{solution}

\begin{solution}{exo:b2:magmatic-differentiation:10}
(a) $z = (0.2 + 1.5 - 1)/(0.2 - 1) = 0.7/(-0.8) = -0.875$~; $F^{-z} = 0.7^{0.875} =
0.7319$~; $C^*/C_0 = 0.2\times0.5/(0.2 - 1 + 1.5) = 0.1/0.7 = 0.1429$, donc $C/C_0 =
0.7319 + 0.1429\times(1 - 0.7319) = 0.7319 + 0.0383 = 0.770$. (b) Rayleigh~:
$0.7^{0.5} = 0.837$. Le \omterm{def:b2:magma-emplacement:magma}{magma} en AFC a moins de strontium pour deux raisons~: pour perdre
\qty{30}{\percent} de sa masse tout en gagnant $r$ fois ce qu'il cristallise, il doit cristalliser
$(1 - 0.7)/(1 - 0.2) = 0.375$ au lieu de 0.30, et le strontium qu'il gagne
provient d'une roche deux fois moins riche. (c) $(0.7319\times0.7035 + 0.0383\times0.7200)/0.7702
= (0.5149 + 0.0276)/0.7702 = 0.70432$~: le rapport a monté de 0.0008
alors que la seule cristallisation l'aurait laissé à 0.7035.
\end{solution}

\begin{solution}{exo:b2:magmatic-differentiation:11}
(a) $1.1\times(1100 - 200) + 250 = 990 + 250 = \qty{1240}{J}$ par gramme. (b)
$25.4\times1240 = \qty{31496}{J}$, proche des \qty{31350}{J} de Wilcox. (c)
$31496/370 = \qty{85.1}{g}$ de plagioclase. (d) Le bilan de masse n'a retiré que $2.9 +
9.6 = \qty{12.5}{g}$ de cristaux, dont la chaleur latente (environ \qty{4600}{J}) couvre moins
d'un sixième du besoin~; $r = 25.4/12.5 = 2.03$, soit deux fois plus de roche assimilée
que de cristaux formés, alors que le bilan thermique admet un $r$ bien inférieur à 1.
Le \omterm{def:b2:magma-emplacement:magma}{magma} éruptif ne peut pas avoir fourni lui-même la chaleur~: un volume de \omterm{def:b2:magma-emplacement:magma}{magma} plus grand
que celui qui a été émis a dû cristalliser et se refroidir en profondeur, hors de vue.
\end{solution}

\begin{solution}{exo:b2:magmatic-differentiation:12}
(a) Étape 1, $F$ de 1 à 0.8 avec $D = 5$~: facteur $0.8^{4} = 0.4096$. Étape 2,
de 0.8 à 0.5 avec $D = 1.5$~: la loi s'applique au liquide restant, donc le
facteur vaut $(0.5/0.8)^{0.5} = 0.7906$. Nickel~: $250\times0.4096\times0.7906 =
\qty{81.0}{ppm}$. (b) $250\times0.5^{4} = \qty{15.6}{ppm}$, cinq fois trop bas. (c)
La loi de Rayleigh suppose un $D$ constant~; quand l'assemblage cristallisant change,
$D$ saute, et la loi doit être repartie de la composition du liquide atteinte au
changement, avec $F$ mesuré par rapport à ce liquide.
\end{solution}

\begin{solution}{pb:b2:magmatic-differentiation:1}
\noindent\textbf{Partie I.}

\textbf{1.} $\bar D_{\ce{Ni}} = 0.20\times5.0 + 0.50\times0 + 0.30\times2.0 = 1.60$.

\textbf{2.} $\bar D_{\ce{Zr}} = 0.20\times0.0007 + 0.50\times0.001 + 0.30\times0.195 =
0.00014 + 0.0005 + 0.0585 = 0.059$.

\textbf{3.} Zone supérieure~: $\bar D_{\ce{Ni}} = 0.45\times2.0 = 0.90$~; $\bar D_{\ce{Zr}} =
0.55\times0.001 + 0.45\times0.195 = 0.088$.

\textbf{4.} Le nickel est compatible dans la zone inférieure (1.60) et devient légèrement
incompatible (0.90) quand l'olivine disparaît~; le zirconium est incompatible dans les deux.

\textbf{5.} Le clinopyroxène~: il fournit $0.0585$ des $0.059$ de la zone inférieure, presque
tout, car c'est le seul minéral dont le coefficient s'éloigne nettement de zéro.

\textbf{6.} $(C_{\ce{Ni}}/C_{\ce{Zr}})$ varie de $F^{D_{\ce{Ni}} - D_{\ce{Zr}}}$~;
par division par deux, $0.5^{1.60 - 0.059} = 0.5^{1.541} = 0.344$.

\noindent\textbf{Partie II.}

\textbf{7.} $1.60\times150 = \qty{240}{ppm}$.

\textbf{8.} Le liquide contenait $190/1.60 = \qty{118.75}{ppm}$~; $F = (118.75/150)^{1/0.6} =
0.7917^{1.667} = 0.68$.

\textbf{9.} $80\times0.68^{-0.941} = 80\times1.438 = \qty{115}{ppm}$.

\textbf{10.} Fractionnée~: $150\times0.68^{0.6} = 150\times0.793 = \qty{119.0}{ppm}$~;
à l'équilibre~: $150/(1.6 + 0.68\times(1 - 1.6)) = 150\times0.839 = \qty{125.8}{ppm}$.
La différence n'est que de \qty{6}{\percent} lorsqu'un tiers du \omterm{def:b2:magma-emplacement:magma}{magma} a cristallisé~:
pour un élément modérément compatible, les deux lois restent proches au début.

\textbf{11.} $150\times0.55^{0.6} = \qty{104.8}{ppm}$.

\textbf{12.} Juste en dessous~: $1.60\times104.8 = \qty{167.7}{ppm}$~; juste au-dessus~:
$0.90\times104.8 = \qty{94.3}{ppm}$. Le profil de nickel des \omterm{def:b2:magmatic-differentiation:cumulate}{cumulats} chute
brutalement de \qty{44}{\percent} au niveau où l'olivine disparaît des
couches~: une marche, et non une courbe régulière.

\textbf{13.} Liquide~: $104.8\times(0.40/0.55)^{0.90 - 1} = 104.8\times1.032 =
\qty{108.2}{ppm}$~; \omterm{def:b2:magmatic-differentiation:cumulate}{cumulat} $0.90\times108.2 = \qty{97.4}{ppm}$. Avec $\bar D < 1$
le nickel augmente désormais lentement dans le liquide et dans les \omterm{def:b2:magmatic-differentiation:cumulate}{cumulats}~: le profil, après sa
marche, remonte doucement.

\noindent\textbf{Partie III.}

\textbf{14.} $x_{\mathrm{ol}}(m_{\mathrm{ol},j} - d_j) + x_{\mathrm{pl}}
(m_{\mathrm{pl},j} - d_j) + x_{\mathrm{cpx}}(m_{\mathrm{cpx},j} - d_j) = 100\,(p_j -
d_j)$.

\textbf{15.} \ce{MgO}~: $37.16\,x_{\mathrm{ol}} - 4.9\,x_{\mathrm{pl}} +
11.1\,x_{\mathrm{cpx}} = 410$~; \ce{CaO}~: $-7.3\,x_{\mathrm{ol}} +
7.06\,x_{\mathrm{pl}} + 13.2\,x_{\mathrm{cpx}} = 420$~; \ce{Al2O3}~:
$-14.1\,x_{\mathrm{ol}} + 17.6\,x_{\mathrm{pl}} - 11.1\,x_{\mathrm{cpx}} = 190$.

\textbf{16.} En résolvant (par élimination ou par la règle de Cramer)~: $x_{\mathrm{ol}} = 9.03$,
$x_{\mathrm{pl}} = 30.84$, $x_{\mathrm{cpx}} = 20.32$ grammes pour \qty{100}{g}.

\textbf{17.} Cristaux $9.03 + 30.84 + 20.32 = \qty{60.19}{g}$~; liquide
\qty{39.81}{g}~: $F = 0.398$.

\textbf{18.} $(4950 - 9.03\times39.19 - 30.84\times50.54 - 20.32\times51.5)/39.81 =
(4950 - 353.9 - 1558.7 - 1046.5)/39.81 = 1990.9/39.81 = \qty{50.0}{\percent}$, comme
dans le \omterm{def:b2:magma-emplacement:bodies}{dyke}~: l'oxyde non utilisé comme témoin est reproduit.

\textbf{19.} Les cristaux retirés sont $9.03/60.19 = \qty{15}{\percent}$ d'olivine,
\qty{51}{\percent} de plagioclase et \qty{34}{\percent} de clinopyroxène, entre les
proportions des deux zones. Ils correspondent à l'assemblage inférieur qui cristallise de $F = 1$ à
0.55 ($0.45\times20 = 9.0$, $0.45\times50 = 22.5$, $0.45\times30 = 13.5$), suivi de l'assemblage supérieur de 0.55 à
0.40 ($0.15\times55 = 8.25$, $0.15\times45 = 6.75$)~: 9.0,
30.75 et 20.25 grammes.

\textbf{20.} $0.30/0.75 = 0.40$.

\noindent\textbf{Partie IV.}

\textbf{21.} $0.55^{0.059 - 1} = 0.55^{-0.941} = 1.755$; $80\times1.755 =
\qty{140.4}{ppm}$.

\textbf{22.} $188/140.4 = 1.339 = (F/0.55)^{0.088 - 1}$, donc $F = 0.55\times1.339^{-1/
0.912} = 0.55\times0.726 = 0.399$.

\textbf{23.} $(188/80)^{1/(0.07 - 1)} = 2.35^{-1.075} = 0.399$. Pour les \omterm{def:b2:elements-in-the-earth:partition}{éléments incompatibles}
$C_L/C_0 \approx 1/F$ quelle que soit la petite valeur de $D$~: la réponse dépend à peine
des détails de $D$, ce qui fait des \omterm{def:b2:elements-in-the-earth:partition}{éléments incompatibles} les meilleures jauges de $F$.

\textbf{24.} Le plagioclase, qui ne contient pas de fer, forme la moitié des cristaux, et le
pyroxène et l'olivine réunis contiennent moins de fer qu'ils ne contiennent de magnésie~: le liquide
s'enrichit en fer et s'appauvrit en magnésie sans gagner de silice. Cet enrichissement en fer
à silice presque constante est la marque de la série tholéiitique du
chapitre suivant.

\textbf{25.} Des cristaux cumulus purs contiendraient $0.088\times188 = \qty{16.5}{ppm}$. Avec une
fraction massique $t$ de liquide piégé, $16.5(1 - t) + 188\,t = 47$, donc $t = (47 -
16.5)/(188 - 16.5) = 30.5/171.5 = 0.18$~: environ \qty{18}{\percent} de liquide piégé.

\textbf{26.} Les trois estimations, par les éléments majeurs (0.398), par la
potasse (0.40) et par le zirconium (0.399), concordent~: quand la couche la
plus haute a cristallisé, il restait encore liquide
\textbf{\qty[text-series-to-math=true]{40}{\percent}} du \omterm{def:b2:magma-emplacement:magma}{magma} d'origine.
\end{solution}
